\phantomsection
\chapter*{Pendakian Terindah}
\addcontentsline{toc}{section}{Pendakian Terindah --- Hadjar Azriel Miftahul Khoir} 
\penulis{Hadjar Azriel Miftahul Khoir}

Pada hari minggu tanggal 16 Agustus kemarin aku bersama temanku melakukan perjalanan ke gunung ciremai. Kami melakukan perjalanan dengan menggunakan kendaraan bermotor. Dan kami melewati beberapa kota, Kami rehat sejenak di kota sumedang sambil mencoba makanan khas kota sumedang yaitu tahu sumedang. Lalu kita melanjutkan perjalanan buat ke basecamp ciremai, setelah sampai basecamp keadaan langit masih gelap, akupun berdiri di depan sebuah warung kecil di kaki Gunung Ciremai. Udara pagi terasa begitu dingin, ku gosokan kedua tanganku. Tidak lupa aku periksa kembali tas perbekalan agar tidak ada kebutuhan yang tertinggal. Dimulai dari hal -hal terkecil seperti air, jas hujan, makanan ringan, kotak P3K, dan kebutuhan logistik lainnya.


Aku Bersama temanku memastikan segala sesuatunya dan semua sudah siap sedia, karena ini pertama kalinya kita mendaki ke Gunung Ciremai.


Dihadapan kami, Gunung Ciremai berdiri dalam kegelapan. Langit malam yang tertutup karena kabut sehingga belum terlihat puncaknya. Meski begitu, membuat jantungku berdebar membayangkan perjalanan yang akan kami lalui.


Sudah sejak lama aku ingin sekali mendaki Gunung Tertinggi di jawa Barat itu. Tetapi keinginan tersebut selalu tertunda karena kesibukan sekolah dan berbagai alasan lainnya. Pada akhirnya kami berangkat pada tanggal 16 Agustus karena besok nya adalah tanggal 17 Agustus yang bertepatan dengan Hari Kemerdekaan Republik Indonesia, sekaligus kami pun bisa mengibarkan bendera di puncak gunung ciremai, betapa indahnya dan betapa bangganya kami sebagai bangsa Indonesia dan pemuda harapan bangsa bisa melakukan upacara di puncak Gunung Ciremai.


Kalau bukan sekarang lagi, kapan lagi dan dengan adanya momentum acara tersebut, yang membuatku akhirnya berani mencoba melakukan pendakian ke Puncak Gunung Ciremai. Dengan Semangat juang 45 dan api keberanian yang berkobar dalam dada.


Setelah memastikan semua perlengkapan aman, kami mulai berjalan. Jalur pendakian perlahan membawa kami memasuki hutan. Suara langkah kaki menjadi satu-satunya suara yang terdengar. Sesekali terdengar suara kicau burung dari kejauhan.


Pada awal perjalanan, aku masih merasa sangat bersemangat. Kakiku terasa ringan dan napasku teratur. Aku bahkan beberapa kali mengajak temanku bercanda.


Semakin jauh kami berjalan, jalur semakin menanjak. Tanah yang kami pijak mulai terasa licin. Beberapa kali aku harus berpegangan pada akar pohon agar tidak terpeleset.


Sekitar dua jam perjalanan, napasku mulai terengah-engah.


“Berhenti dulu,” kata temanku.


Aku langsung meletakkan tasku. Dan bertanya apakah kamu capek?


Temanku mengangguk.


Kami duduk di sebuah tempat yang cukup datar. Dari sela-sela pepohonan, cahaya matahari mulai masuk. Pagi perlahan menggantikan malam. Aku membuka botol air dan meminumnya sedikit. Temanku mengambil biskuit dari tasnya lalu membaginya menjadi dua.


“Setengah buat kamu,” katanya.


“Kenapa cuma setengah?”


“Biar adil.”


“Peliiit.”


Kami tertawa.


Perjalanan kemudian kami lanjutkan. Namun, rasa lelah semakin terasa. Bahuku mulai sakit karena beban tas, sedangkan kakiku terasa berat. Dalam hati aku mulai bertanya-tanya apakah aku sanggup mencapai puncak.


Beberapa kali aku berjalan lebih dulu. Namun setiap kali menyadari temanku tertinggal, akupun berhenti dan menunggu.


“Pelan-pelan saja,” kataku. “Gunungnya enggak akan lari.”


Temanku tersenyum mendengar perkataanku.


Menjelang sore, kami sampai di tempat yang cukup terbuka untuk mendirikan tenda. Kabut mulai turun. Udara terasa jauh lebih dingin dibandingkan saat kami memulai pendakian.


Kami segera mendirikan tenda sebelum hari benar-benar gelap.


Setelah tenda berdiri, kami memasak mi instan dan membuat minuman hangat. Sederhana, tetapi rasanya sangat nikmat.


Kami duduk di depan tenda sambil memandangi langit.


“Lihat,” kata temanku sambil menunjuk ke atas.


Aku mendongak.


Bintang-bintang terlihat begitu jelas. Tidak ada lampu kota yang mengganggu. Langit seakan menjadi lebih luas dan dekat.


“Indah banget,” kataku pelan.


Makanya orang-orang rela capek-capek naik gunung.”


Malam itu kami berbicara banyak hal. Tentang sekolah, keluarga, cita-cita, dan kehidupan setelah lulus. Sesekali kami tertawa mengingat kejadian-kejadian lucu yang pernah kami alami.


Namun, malam semakin larut. Kami akhirnya masuk ke dalam tenda dan beristirahat.


Pukul dua dini hari, Aku membangunkan temanku.


“Bangun.”


Aku membuka mata dengan malas.


“Sudah waktunya.”


Aku melihat jam tangan. Benar saja, waktunya sudah menunjukkan pukul dua.


Kami bersiap dengan cepat. Jaket, sarung tangan, dan senter sudah terpasang. Setelah memastikan tenda aman, kami mulai berjalan menuju puncak.


Jalur malam terasa berbeda. Gelap membuat setiap langkah harus lebih hati-hati. Cahaya senter menjadi satu-satunya penerangan.


Tidak banyak percakapan di antara kami. Kami lebih banyak diam dan berkonsentrasi menjaga langkah.


Semakin mendekati puncak, medan semakin berat. Angin bertiup kencang. Batu-batu besar membuat kami harus menggunakan tangan untuk membantu menarik tubuh.


Aku mulai merasa sangat lelah.


“Kita istirahat sebentar.”


Kami duduk di sebuah batu. Aku terdiam. kalau cuma karena capek, kita istirahat dulu. Setelah itu coba lagi. Sedikit demi sedikit.”


Aku memandang jalur di depan kami. Puncak masih belum terlihat. Namun, entah kenapa, semangatku kembali.


Kami melanjutkan perjalanan.


Satu langkah.


Berhenti.


Satu langkah lagi.


Berhenti.


Begitulah kami berjalan. Tidak cepat, tetapi terus bergerak.


Sekitar pukul lima pagi, langit di sebelah timur mulai berubah warna. Gelap perlahan menjadi biru tua, kemudian jingga.


“Sebentar lagi”


Aku mengangkat kepala.


Di depan kami, jalur menuju puncak mulai terlihat.


Aku mengumpulkan sisa tenaga dan terus berjalan.


Beberapa menit kemudian, kakiku akhirnya menginjak tanah datar di puncak.


Kami berhasil.


Aku berdiri diam sambil menatap sekeliling. Rasa lelah yang sejak tadi memenuhi tubuh seolah hilang untuk beberapa saat.


Matahari mulai muncul dari balik awan. Langit berubah menjadi warna keemasan. Gunung-gunung di kejauhan terlihat seperti pulau-pulau yang mengambang di atas lautan awan.


Aku tidak bisa berkata apa-apa.


Temanku berdiri di sampingku.


Tidak lupa aku keluarkan bendera merah putih dalam tasku, dan kami pun mengibarkan berndera merah putih di puncak gunung ciremai.


Ditengah angin yang  bertiup sangsaka Merah Putih terus berkibar, seolah menggambarkan semangat keberanian dan perjuangan.


Kami kemudian duduk berdampingan. Tidak ada percakapan selama beberapa menit. Kami hanya menikmati pemandangan.


Saat itu aku sadar, ternyata tujuan pendakian bukan hanya tentang mencapai puncak.


Perjalanan menuju puncaklah yang membuat semuanya berarti.


Aku belajar bahwa rasa lelah tidak selalu berarti harus berhenti. Kadang kita hanya membutuhkan waktu untuk beristirahat sebelum kembali melangkah. Aku juga belajar bahwa seorang teman bukan hanya seseorang yang berjalan bersama kita ketika semuanya mudah. Teman adalah orang yang tetap berada di samping kita ketika langkah mulai terasa berat.


Namun, pengalaman sederhana itu akan menjadi salah satu kenangan yang tidak akan pernah kulupakan.


Sebelum turun, aku mengambil satu foto bersama temanku dengan latar belakang lautan awan. dan Bendera Merah Putih.

 
“Foto ini jangan sampai hilang,” kataku.


“Tenang. Kalau hilang, kita naik lagi buat foto ulang.”


Aku tertawa.


“Jangan bercanda.”


“Serius.”


Kami tertawa bersama.


Kemudian kami mulai turun meninggalkan puncak Ciremai.


Kali ini langkah kami terasa lebih ringan.


Bukan karena perjalanan menjadi lebih mudah, tetapi karena kami tahu bahwa kami telah berhasil melewati sesuatu yang sebelumnya kami pikir tidak mampu kami lakukan.


Gunung Ciremai mengajarkan kami satu hal sederhana: puncak memang indah, tetapi keberanian untuk terus melangkah jauh lebih berharga.


Dan sejak hari itu, setiap kali aku menghadapi sesuatu yang terasa berat dalam hidup, aku selalu mengingat perjalanan tersebut.


Satu langkah.


Kemudian satu langkah lagi.


Sampai akhirnya kita tiba di tempat yang selama ini kita tuju.


Seolah semangat kemerdekaan terlihat sangat nampak setelah melakukan perjalanan pendakian menuju puncak gunung ciremai yang merupakan sebuah perjuangan.


Dan mengingatkan kita akan perjuangan dan pengorbanan para pahlawan dalam memperjuangkan kemerdekaan Negara kita Republik Indonesia.


Walaupun kita tidak secara langsung berperang, namun dengan melakukan perjalanan pendakian ke gunung ciremai kali ini, menimbulkan semangat juang 45. Sebagai generasi penerus bangsa sudah seharusnya kita tetap menjaga semangat kemerdekaan yang telah di peroleh dengan cara melakukan hal yang positif dan bermanfaat.


Semoga perjalanan pendakian ke Gunung Ciremai kali ini bisa bermanfaat dan menjadikan motivasi kepada teman teman yang lain. 

